\section{Certificates and query-specific evidence costs}
\label{sec:post-cstat-theory}

The object of inference is a fixed-design moment. The probability model permits
heterogeneous row means and one connected dependency graph; it does not require
independent repeated macro-blocks or equal block targets. We condition throughout
on a legitimate exogenous field $\mathcal D$, augmented only by training and
discovery information whose random sources are independent of certification.
A schema can check a declared source construction, but cannot establish the
absence of unrecorded common causes from a single realized log.

\paragraph{Signed moments and support.}
For coordinate $\theta$, let $c_\theta$ identify its response channel, and write
$Y_{u,\theta}$, $\mu_{us,\theta}$ and $\widehat\mu_{us,\theta}$ for the response,
identified exposure-specific mean and fixed nuisance fit. The finite exposure
law $e_{us}=\Pr(S_u=s\mid\mathcal D)$ is known exactly. Define
\begin{align}
 M_\theta
 &=a_\theta+\frac{1}{B_\theta}
     \sum_{u=1}^N\sum_s d_{us,\theta}\mu_{us,\theta},
 &B_\theta&=N\,\mathrm{normalizer}_\theta>0,
 \label{eq:post-target}\\
 D_\theta
 &=a_\theta+\frac{1}{B_\theta}
     \sum_{u,s}d_{us,\theta}\widehat\mu_{us,\theta},
 &R_{u,\theta}
 &=\frac{d_{uS_u,\theta}}{e_{uS_u}}
       (Y_{u,\theta}-\widehat\mu_{uS_u,\theta}).
 \label{eq:post-residual}
\end{align}
Here $a_\theta$, $B_\theta$ and the signed coefficients $d_{us,\theta}$ are
$\mathcal D$-measurable. A level query uses $d=hq$; a paired policy contrast
uses $d=h(q-q_0)$. A subgroup ratio uses its known positive design mass in
$\mathrm{normalizer}_\theta$; a zero-mass ratio is undefined. Every nonzero
coefficient requires $e_{us}>0$. Exact exposure probabilities identify the
counterfactual mean only together with the stated exposure sufficiency and
assignment assumptions. A policy interpretation additionally requires that
$q$ arise from a compatible joint assignment law.

Support is a property of the question. For bounded responses with exposure
law $e=(0,1/2,1/2)$, consider
$q_0=(1/5,3/10,1/2)$ and $q_1=(1/5,1/2,3/10)$.
Neither policy level is point identified when the first state's mean is
unknown. Their signed contrast has coefficient
$q_1-q_0=(0,1/5,-1/5)$ and uses only observed states. The common unsupported
contribution cancels. Every matched contrast-aware method receives this
same supported query.

\begin{assumption}[Conditional graph and residual envelopes]
\label{ass:post-core}
The desired exposure-specific means satisfy
$\mathbb E[Y_{u,\theta}\mid S_u=s,\mathcal D]=\mu_{us,\theta}$.
The entire row score/residual vector has a sufficient conditional dependency
graph $G$: disjoint vertex sets with no crossing edge are conditionally
independent given $\mathcal D$. For a declared $1<p\leq2$, known numbers
$b_{us,\theta}\geq0$ satisfy
\begin{equation}
 \mathbb E\!\left[
 |Y_{u,\theta}-\widehat\mu_{us,\theta}|^p
 \mid S_u=s,\mathcal D\right]\leq b_{us,\theta}^p.
 \label{eq:post-channel-envelope}
\end{equation}
The query family, all confidence allocations and all tuning choices are fixed
conditional on $\mathcal D$ before certification outcomes are examined.
\end{assumption}

A sufficient constructive graph connects every pair of rows sharing any
independent primitive source used in logging or in any retained response
channel. Its sufficiency follows from disjoint source collections, not from
observed correlations or a domain network alone. Let $A_\theta$ be the rows
whose residual is not identically zero by design. Use the induced graph
$G[A_\theta]$ and the smallest number $\chi_\theta$ among the available,
verified proper colorings. This is an available coloring count, rather than
a claim to compute the chromatic number. All source-induced edges between
retained rows remain present even when other rows are removed.

Set
\begin{equation}
 v_{u,\theta}
 =\sum_{s:e_{us}>0}|d_{us,\theta}|^p e_{us}^{1-p}b_{us,\theta}^p,
 \qquad
 c_p=\frac{(p-1)^{p-1}}{p^p}.
 \label{eq:post-v}
\end{equation}
These bound raw residual moments; a centered-noise moment cannot be substituted
without accounting for nuisance error. In bounded responses $Y\in[L,U]$, the
legal residual magnitude is $\max\{|L-\widehat\mu|,|U-\widehat\mu|\}$.
In particular, midpoint nuisance on $[0,1]$ gives magnitude $1/2$ for every
eligible comparison method.

\begin{theorem}[Heterogeneous conditional-graph residual certificate]
\label{thm:post-core}
Under Assumption~\ref{ass:post-core}, choose positive $\mathcal D$-measurable
thresholds $\tau_{u,\theta}$ on the active support and let
$T_\tau(x)=\operatorname{sign}(x)\min\{|x|,\tau\}$. For a coordinate budget
$\delta_\theta\in(0,1)$, put $t_\theta=\log(2/\delta_\theta)$ and define
\begin{align}
 \widehat M_\theta
 &=D_\theta+\frac{1}{B_\theta}
           \sum_{u\in A_\theta}T_{\tau_{u,\theta}}(R_{u,\theta}),
 \label{eq:post-estimator}\\
 r_\theta
 &=\frac{1}{B_\theta}\left[
 c_p\sum_{u\in A_\theta}v_{u,\theta}\tau_{u,\theta}^{1-p}
 +\sqrt{2\chi_\theta t_\theta
           \sum_{u\in A_\theta}v_{u,\theta}\tau_{u,\theta}^{2-p}}
 +\frac{2}{3}\chi_\theta t_\theta
                   \max_{u\in A_\theta}\tau_{u,\theta}\right].
 \label{eq:post-radius}
\end{align}
Then $\Pr(|\widehat M_\theta-M_\theta|>r_\theta\mid\mathcal D)
\leq\delta_\theta$. For a finite frozen family $\Theta$ with
$\sum_{\theta\in\Theta}\delta_\theta\leq\delta$, all coordinate intervals
cover simultaneously with conditional probability at least $1-\delta$.
If no random residual remains, its radius is zero and the known direct term
is retained.
\end{theorem}

The three terms in~\eqref{eq:post-radius} pay for clipping bias, graph-adjusted
fluctuation and bounded range, respectively. The target is the original
unclipped moment. Clipping without the first term targets a different object
unless a separate argument removes the bias. One threshold per query is the
special case $\tau_{u,\theta}=\tau_\theta$ used by the initial implementation.
The proof, including the graph MGF step, appears in
Appendix~\ref{app:post-proofs}.
Each family statement concerns one predeclared construction per coordinate,
including a design-only dispatch when used. It does not automatically give
joint coverage over all plotted methods and both nominal levels. Such a claim
would require allocating the total budget across method, level and query.

\begin{corollary}[Locality, design-only tuning and paired scaling]
\label{cor:post-locality}
Removing design-inactive rows leaves the estimand and direct term unchanged,
and permits the induced support graph in Theorem~\ref{thm:post-core}.
Thresholds and eligible certificate engines may be selected by minimizing a
proved radius over a predeclared design-only grid without an additional
outcome-search budget. For a signed contrast scaled by $\eta>0$, scaling its
thresholds by $\eta$ scales its signal and every radius component by $\eta$;
$\eta=0$ has zero random uncertainty. Small policy steps therefore do not by
themselves improve sign-certification power.
\end{corollary}

\paragraph{Sharp bounded comparisons.}
For known response bounds, let $[\ell_{u,\theta},u_{u,\theta}]$ be the actual
possible residual interval, taking extrema over every supported exposure, and
write $w_{u,\theta}=u_{u,\theta}-\ell_{u,\theta}$. With the same query, nuisance,
support, graph and coordinate budget, the untruncated estimator
$\widehat M^{\mathrm{mean}}_\theta=D_\theta+B_\theta^{-1}\sum_u R_{u,\theta}$
has graph Hoeffding and Bernstein radii
\begin{align}
 r_{H,\theta}
 &=\frac{1}{B_\theta}
       \sqrt{\frac{\chi_\theta t_\theta}{2}\sum_u w_{u,\theta}^2},
 \label{eq:post-hoeffding}\\
 r_{B,\theta}
 &=\frac{1}{B_\theta}\left[
       \sqrt{2\chi_\theta t_\theta\sum_u V_{u,\theta}}
       +\frac{\chi_\theta t_\theta}{3}\max_u w_{u,\theta}\right],
 \quad
 V_{u,\theta}=\min\left\{v^{(2)}_{u,\theta},
                               \frac{w_{u,\theta}^2}{4}\right\},
 \label{eq:post-bernstein}
\end{align}
where $v^{(2)}$ is the legal raw second-moment bound from the known bounded
response class. The same bounds and envelope information are supplied to all
matched methods. A generic bias-corrected truncated graph mean with the same
inputs and tuning is algebraically identical to the corresponding
configuration of~\eqref{eq:post-estimator}--\eqref{eq:post-radius}; equality is
a comparison result, not evidence of estimator superiority.

For the historical A design alone, $q=e=(1/2,1/2)$, midpoint nuisance and
reference $1/4$ give score $Y-1/4$. At $N=8192$, $m=1$, $\chi=1$ and
$\delta=.05$,
\begin{equation}
 r_H=\sqrt{\frac{\log 40}{2\cdot8192}}
       =0.015005043614373761,
 \qquad
 .0375>2r_H.
 \label{eq:post-A-fixture}
\end{equation}
This is a deterministic analytic comparison: on its coverage event the lower
bound for target $.2375$ exceeds threshold $.2$. It is not new measured power
and is not included in the old seventeen-test confirmation family.

\begin{proposition}[A scoped subgroup--overlap--dependence obstruction]
\label{prop:post-lower}
Let $g=N\pi$ be a positive integer divisible by $\chi$, set $n=g/\chi$, and
fix $\epsilon\in(0,1]$, $\nu>0$, $1<p\leq2$ and a mean-deviation cap $H>0$.
There is a supported, exactly known logging experiment on $n$ independent
$\chi$-row cliques, with exposure probability $\epsilon$ and unweighted
potential-response residual moment
$\mathbb E|Y(1)-b|^p\leq\nu^p$, for which every estimator of the normalized
subgroup mean has worst-case expected absolute error at least
\begin{equation}
 \frac12\min\left\{H,
 \nu\min\left(1,(2n\epsilon)^{-\frac{p-1}{p}}\right)\right\}.
 \label{eq:post-lower}
\end{equation}
The corresponding raw importance-weighted residual moment is at most
$\epsilon^{1-p}\nu^p$; $\nu$ is the response-residual scale, before
importance weighting. For the raw subgroup moment the bound is multiplied
by $\pi$. In its
unsaturated regime this gives normalized scale
$\nu(\chi/(N\pi\epsilon))^{(p-1)/p}$, up to the displayed constant.
It does not establish graph-by-graph sharpness or matching simultaneous
confidence logarithms. The construction and a non-sharp confidence-width
consequence are in Appendix~\ref{app:post-lower}.
\end{proposition}

A symmetric certificate with radius $r$ certifies a true threshold violation
on its coverage event whenever its true gap satisfies $\gamma>2r$.
This is a sufficient procedure-specific decision condition. The lower bound
is a separate necessary expected-risk obstruction on its stated subclass;
neither statement turns a conservative radius map into a sharp information
limit for every design.

\begin{corollary}[Independent discovery and certified finite selection]
\label{cor:post-selection}
A finite family selected using independent discovery sources may be included
in $\mathcal D$. The simultaneous family must contain every risk coordinate
and utility contrast actually used by selection. On its coverage event,
returning only a candidate whose risk upper bounds meet the task constraints,
whose exact design service/capacity meet their bounds, and whose utility lower
contrast meets the declared usefulness requirement is safe and useful in that
fixed-design sense. In particular,
$\Pr(\mathrm{return}\ \mathrm{and}\ \mathrm{unsafe})\leq\delta$.
There is no corresponding automatic $\delta$ guarantee for
$\Pr(\mathrm{unsafe}\mid\mathrm{return})$ and no theorem here guaranteeing
a positive return rate.
\end{corollary}

\begin{corollary}[Affine mixtures on the simultaneous event]
\label{cor:post-mixtures}
Let $Q_j$ be compatible base joint policies with fixed auditors and normalizers.
For nonnegative weights summing to one, $Q_w=\sum_jw_jQ_j$ is a valid joint
mixture. Simultaneous base intervals propagate to
$[\sum_jw_jL_j,\sum_jw_jU_j]$ for each affine moment, even for weights selected
after seeing the certificates. This is linear propagation of certified base
quantities, not permission to recompute a nonlinear clipped estimator at an
outcome-selected mixture. Nonlinear metrics or moving denominators require
separate analysis.
\end{corollary}

All statements above concern real arithmetic under declared assumptions.
Implementation uses a separate conservative numerical rule and must preserve
empty confidence intersections, valid but vacuous intervals, unsupported
queries and numerical failures as distinct recorded outcomes. No adaptive
history, arbitrary stopping-time or future-deployment guarantee is asserted.
